\documentclass[11pt]{article}

\input{Shared.tex}

\parindent0in
\pagestyle{plain}
\thispagestyle{plain}

\usepackage{csquotes}
\usepackage[shortlabels]{enumitem}

\newcommand{\dated}{\today}
\newcommand{\token}[1]{\langle \text{#1} \rangle}

\begin{document}

\textbf{Introduction to the Theory of
Computation}\hfill\textbf{\myname}\\[0.01in]
\textbf{Chapter 8: Space Complexity}\hfill\textbf{\dated}\\
\smallskip\hrule\bigskip

\begin{problem}{8.15}
Consider the following two-person version of the language PUZZLE that was described in Problem 7.28. Each player starts with an ordered stack of puzzle cards. The players take turns placing the cards in order in the box and may choose which side faces up. Player I wins if all hole positions are blocked in the final stack, and Player II wins if some hole position remains unblocked. Show that the problem of determining which player has a winning strategy for a given starting configuration of the cards is PSPACE-complete.
\end{problem}

\begin{proof}
Let $2\text{-}PUZZLE = \{\langle C_1, C_2, S, p \rangle \ | \ C_1, \ C_2 \text{ and } S \text{ are ordered stacks of cards,} \\ \text{representing current configuration of the cards of Player 1, Player 2 and stack, such that Player 1} \\ \text{has a winning strategy if player } p \text{ takes the first turn}\}$. To show that $2\text{-}PUZZLE$ is PSPACE-complete, we must show that it is in PSPACE and that all PSPACE problems are polynomial time reducible to it. To prove $2\text{-}PUZZLE \in$ PSPACE, we give a polynomial space algorithm $M$. \\

$M =$ \textquotedblleft On input $\langle C_1, C_2, S, p \rangle$, where $C_1$, $C_2$ and $S$ are ordered stacks of cards, and $p$ is an integer 1 or 2:
\begin{enumerate}
\item If both $C_1$ and $C_2$ are empty, then:
\item \hspace*{0.5cm} If all hole positions are blocked in stack $S$, then \textit{accept}. Otherwise, \textit{reject}.
\item If $p$ = 1, then:
\item \hspace*{0.5cm} If stack $C_1$ is not empty, then:
\item \hspace*{1.0cm} Remove top most card from $C_1$, and place it on top of $S$.
\item \hspace*{1.0cm} Recursively call $M$ on $\langle C_1, C_2, S, 2 \rangle$.
\item \hspace*{1.0cm} Remove top most card from $S$, flip it and place it back on top of $S$.
\item \hspace*{1.0cm} Recursively call $M$ on $\langle C_1, C_2, S, 2 \rangle$.
\item \hspace*{1.0cm} Remove top most card from $S$, flip it and place it on top of $C_1$.
\item \hspace*{1.0cm} If one of the recursive calls accept, \textit{accept}. Otherwise, \textit{reject}.
\item \hspace*{0.5cm} Otherwise, recursively call $M$ on $\langle C_1, C_2, S, 2 \rangle$, and operate as $M$ afterwards.
\item If $p$ = 2, then:
\item \hspace*{0.5cm} If stack $C_2$ is not empty, then:
\item \hspace*{1.0cm} Remove top most card from $C_2$, and place it on top of $S$.
\item \hspace*{1.0cm} Recursively call $M$ on $\langle C_1, C_2, S, 1 \rangle$.
\item \hspace*{1.0cm} Remove top most card from $S$, flip it and place it back on top of $S$.
\item \hspace*{1.0cm} Recursively call $M$ on $\langle C_1, C_2, S, 1 \rangle$.
\item \hspace*{1.0cm} Remove top most card from $S$, flip it and place it on top of $C_2$.
\item \hspace*{1.0cm} If both of the recursive calls accept, \textit{accept}. Otherwise, \textit{reject}.
\item \hspace*{0.5cm} Otherwise, recursively call $M$ on $\langle C_1, C_2, S, 1 \rangle$, and operate as $M$ afterwards.\textquotedblright
\end{enumerate}

To prove that $2\text{-}PUZZLE$ is PSPACE-hard, we give a polynomial space reduction from the special case of $TQBF$ presented in Problem 8.12 to $2\text{-}PUZZLE$\footnote{TQBF restricted to formulas where the part following the quantifiers is in conjunctive normal form is still PSPACE-complete.}. This reduction converts a fully quantified Boolean formula $\phi$ to ordered stacks of cards $C_1$, $C_2$ and $S$, and an integer $p$, such that $\phi$ is true, iff Player 1 has a winning strategy if player $p$ takes the first turn and $C_1$, $C_2$ and $S$ represent the current configuration of the cards of Player 1, Player 2 and stack. Next, we show how to construct $C_1$, $C_2$, $S$ and $p$.
\begin{enumerate}
\item Let $p = 1$.
\item Let $C_1$, $C_2$, and $S$ be empty stacks.
\item Construct a card for each variable in $\phi$ according to the construction given in the solution for Problem 7.28.
\item Card for there-exist bound variable is placed in $C_1$.
\item Card for for-all bound variable is placed in $C_2$.
\item Assign a dummy card with all slots punched to player whose turn is skipped due to repeated there-exist or for-all.
\end{enumerate}
\end{proof}

\end{document}